\documentclass[12pt, a4paper]{article}
\usepackage[utf8]{inputenc}
\usepackage[T1]{fontenc}
\usepackage{geometry}
\usepackage{xcolor}
\usepackage{url}
\usepackage{parskip} % Adiciona espaço entre parágrafos automaticamente

% Configuração das margens
\geometry{
    left=2.5cm,
    right=2.5cm,
    top=2.5cm,
    bottom=2.5cm,
}

% --- Definição de Comandos Personalizados para Agilizar ---
% Comando para iniciar a seção de um revisor
\newcommand{\reviewer}[1]{
    \vspace{1cm}
    \hrule
    \section*{Reviewer \##1}
}

% Comando para a preocupação/comentário do revisor
% Uso: \concern{Número}{Texto do comentário}
\newcommand{\concern}[2]{
    \noindent \textbf{Concern \##1:} \textit{#2}
    \par\vspace{0.3cm}
}

% Comando para a resposta do autor
% Uso: \response{Texto da resposta}
\newcommand{\response}[1]{
    \noindent \textbf{Author response:} #1
    \par\vspace{0.3cm}
}

% Comando para a ação tomada no manuscrito
% Uso: \action{Texto da ação}
\newcommand{\action}[1]{
    \noindent \textbf{Author action:} #1
    \par\vspace{0.5cm}
}
% -----------------------------------------------------------

\begin{document}

% Cabeçalho com Informações do Manuscrito [cite: 1-4]
\noindent \textbf{Original Manuscript ID:} Access-2025-51751 \\
\noindent \textbf{Original Article Title:} ``Performance Indicators for Shape and Position Assessment in Electromagnetic Inverse Scattering'' \\
\noindent \textbf{To:} IEEE Access Editor \\
\noindent \textbf{Re:} Response to reviewers

\vspace{1cm}

% Carta ao Editor [cite: 5-9]
\noindent Dear Editor,

\noindent Thank you for allowing a resubmission of our manuscript, with an opportunity to address the reviewers’ comments.

\noindent We are uploading:
\begin{itemize}
    \item[(a)] our point-by-point response to the comments (below) (response to reviewers, under ``Author’s Response Files''),
    \item[(b)] an updated manuscript with yellow highlighting indicating changes (as ``Highlighted PDF''), and
    \item[(c)] a clean updated manuscript without highlights (``Main Manuscript'').
\end{itemize}

\vspace{0.5cm}

\noindent Best regards,

\noindent Batista \& Adriano.

\newpage

% --- Início das Respostas aos Revisores ---

% Reviewer 1 [cite: 10-18]
\reviewer{1}

\concern{1}{It is not sufficient to only report the values of the proposed metric. For a fair assessment, the authors should evaluate the same dataset using
commonly adopted metrics such as MSE and SSIM, include metrics that assess contrast, such as CNR, and include shape and localization accuracy metrics such as the Dice coefficient, IoU (Intersection over Union), and the Hausdorff distance. A side-by-side comparison on identical data is necessary so that readers can clearly see in which situations the proposed metric is more sensitive or more informative than existing metrics. How useful is the proposed metric compared to conventional metrics? On the same dataset, please compute and report standard metrics such as MSE, SSIM, CNR, Dice, IoU, and Hausdorff distance together with the proposed metric. Analyze how each metric responds to variations
in restoration quality. In particular, if there are cases where traditional metrics fail to distinguish differences in quality but the proposed metric can, please highlight such case studies. These examples would make the advantages of the proposed metric much more convincing.}

\response{
    Thank you for this valuable comment. We agree that the inclusion of additional metrics can enrich the evaluation of the proposed method's performance. However, we would like to clarify the following points:
    
    The metrics CNR, Dice coefficient, IoU (Intersection over Union), and Hausdorff distance have never been applied in the literature on algorithms for solving the electromagnetic inverse scattering problem. Therefore, investigating the best way to apply these metrics would constitute a separate study, which falls outside the scope of the current article.
    
    Nevertheless, we decided to include comparisons with MSE and SSIM metrics, which are already commonly used in the literature on algorithms for general geometries in electromagnetic inverse scattering problems. We added dedicated subsections (Sections 4.1.4 and 4.2.4) presenting side-by-side comparisons on identical datasets, demonstrating how each metric responds to variations in reconstruction quality. These comparisons highlight specific cases where the proposed indicators provide distinct insights compared to traditional metrics, particularly regarding their reduced sensitivity to contrast estimation errors and background medium variations.
}

\action{We updated the manuscript by including subsections 4.A.4 and 4.B.4 with side-by-side comparisons of MSE, SSIM, and the proposed metrics on identical datasets, highlighting their responses to variations in reconstruction quality.}

\concern{2}{Especially in medical imaging, where clinical interpretation is critical, it is important to show how well the proposed metric values correlate with subjective assessments by clinical experts (e.g., physicians, radiologists, or other experienced observers). At present, the manuscript lacks such expert-based or clinically oriented validation. Including expert scoring of image quality or restoration adequacy, and analyzing its correlation with the proposed metric, would substantially strengthen the work. If the value of the proposed metric is high, can we truly say that the image has been accurately reconstructed? Please provide clear theoretical and/or empirical justification for this claim. For example, we recommend including analyses that show the correlation between the proposed metric and expert visual scores, demonstrating that the metric is consistent with clinical or expert judgment.}

\response{We agree that for medical diagnostic tools, clinical validation is the gold standard. However, we would like to clarify the scope of this work. The proposed metric is intended as a physics-based quantitative indicator for the algorithmic performance of image reconstruction, rather than a clinical diagnostic tool at this stage.

Since this study relies on synthetic and laboratory-controlled data where the Ground Truth (exact physical profile) is known mathematically, we argue that validation against this Ground Truth provides an objective verification that supersedes subjective observation.

It is also worth noting that other works in the literature that proposed reconstruction algorithms for electromagnetic inverse scattering problems and performed experiments with the same dataset did not conduct clinical validation or expert assessment, focusing only on evaluating algorithmic performance in a more challenging scenario, as is the case with the Breast Phantom Dataset.}

\action{We updated the manuscript by adding a discussion in the Breast Phantom Study section (Section 4.C) acknowledging that while the proposed indicators are validated against known ground truth in controlled scenarios, future work could explore correlation with clinical expert assessments for real-world medical applications. We also clarified that the indicators are designed as algorithmic performance metrics for algorithm development and benchmarking, rather than as clinical diagnostic tools requiring expert validation at this stage \cite{contzano2023fast,wu2024physics,shah2018fast,batista2025eispy2d}.}

% Reviewer 2
\reviewer{2}

\concern{1}{The data source gives 3D phantoms (full 3D models), but the scattering model is in 2D. It would be useful to briefly explain what limitations arise when a 3D structure is simplified to 2D.}
\response{Thank you for this important observation. We acknowledge that using a 2D cross-section from a 3D phantom introduces certain limitations. The main limitation is that the 2D model neglects out-of-plane scattering effects and the complete three-dimensional electromagnetic interactions present in the actual breast tissue. This simplification can affect the accuracy of the scattered field representation and consequently impact the reconstruction results.

However, it is important to note that this approach of extracting 2D cross-sections from 3D breast phantoms has been previously adopted in related studies in the electromagnetic inverse scattering literature \cite{contzano2023fast,wu2024physics,shah2018fast,batista2025eispy2d}. While it does not capture the complete three-dimensional nature of the problem, this methodology still provides a challenging and realistic scenario for evaluating reconstruction algorithm capabilities, particularly for demonstrating the applicability of performance indicators. The 2D formulation offers computational efficiency while maintaining sufficient complexity to assess algorithm performance in scenarios with realistic tissue distributions and contrasts.}
\action{We updated the manuscript by adding an explanation in the Breast Phantom Study section (Section 4.C) clarifying that while the 2D approximation does not capture the complete three-dimensional electromagnetic interactions and out-of-plane scattering effects, this approach has been previously validated in the literature as a computationally efficient method for evaluating algorithm performance with realistic tissue configurations.}

\concern{2}{When the original objects have a few pixels (i.e., small targets), the shape metric divides by a very small number (i.e., TP), so the error percentage becomes huge and unstable. A different shape metric (e.g., normalized, etc.) could address this issue.}

\response{Thank you for raising this important concern. We acknowledge that when the total number of scatterer pixels (TP) is very small, the percentage-based shape error indicator can indeed produce large values. However, we argue that this behavior is not necessarily a limitation, but rather an inherent characteristic that reflects the actual challenge of evaluating reconstruction quality for very small targets.

First, it is important to clarify the minimum value of TP is 1, since we are not assuming cases where there is no scatterer present in the domain of interest. In these cases, the shape error indicator is well-defined and it is the number of misclassified pixels (FP + FN) that determines the error percentage.

Second, when dealing with very small scatterers (few pixels), even small absolute errors in shape reconstruction represent significant relative errors, which should be reflected in the performance metric. For instance, if a scatterer occupies only 10 pixels and the reconstruction incorrectly classifies 15 pixels, this represents a 150\% error, which accurately indicates that the reconstruction has more misclassified pixels than the actual object size. This is meaningful information about algorithm performance in challenging scenarios.

Third, the position error indicator $\zeta_P$ complements the shape error indicator in such scenarios. For small targets where shape reconstruction is particularly challenging, the position indicator provides a more stable performance measure, as demonstrated in our experiments in Section 4.B (Position Detection Study), which specifically addresses scenarios with small, high-contrast scatterers.

Finally, we note that the benchmarking studies presented in the manuscript (Sections 4.A.3 and 4.B.3) successfully use these indicators to compare algorithm performance across 30 different instances with varying geometries, demonstrating their practical utility and stability in statistical analyses.}

\action{We updated the manuscript by adding a clarifying note in Section 3.A (Shape Error) explaining that: (1) the minimum value of TP is constrained to 1 to ensure numerical stability; (2) large percentage values for small targets meaningfully reflect the relative difficulty of shape reconstruction in such scenarios; and (3) the position error indicator provides a complementary and more stable measure for small target scenarios, as demonstrated in the Position Detection Study.}


% Notas Finais

\vspace{2cm}
\hrule
\vspace{0.5cm}
\noindent \small{\textbf{Note:} References suggested by reviewers should only be added if it is relevant to the article and makes it more complete. Excessive cases of recommending non-relevant articles should be reported to \url{ieeeaccesseic@ieee.org}.}

% Bibliografia
\bibliographystyle{IEEEtran}
\bibliography{IEEEabrv,../paper/mybib.bib}

\end{document}